\documentclass[10.5pt,a4paper,UTF8]{ctexart}

\usepackage[left=2.05cm,right=2.05cm,top=2.0cm,bottom=2.0cm]{geometry}
\usepackage{amsmath,amssymb}
\usepackage{booktabs,longtable,array,tabularx}
\usepackage{xcolor}
\usepackage{hyperref}
\usepackage{xurl}
\usepackage{fancyhdr}
\usepackage{titlesec}
\usepackage{enumitem}
\usepackage{listings}
\usepackage{fontspec}
\usepackage{xeCJK}

\setmainfont{Microsoft YaHei}
\setCJKmainfont{Microsoft YaHei}
\setmonofont{Consolas}

\definecolor{mainblue}{RGB}{18,82,128}
\definecolor{softblue}{RGB}{232,242,250}
\definecolor{softgreen}{RGB}{235,247,241}
\definecolor{darkgreen}{RGB}{30,115,82}
\definecolor{codebg}{RGB}{248,248,248}
\definecolor{rulegray}{RGB}{210,215,220}

\hypersetup{
  colorlinks=true,
  linkcolor=mainblue,
  urlcolor=mainblue,
  pdftitle={基于 STM32F411 与 MAX30102 的脉搏测试仪验收说明},
  pdfauthor={项目组},
  pdfsubject={项目验收说明书},
  pdfkeywords={STM32F411, MAX30102, PPG, AF, HRV, Stress}
}

\pagestyle{fancy}
\fancyhf{}
\fancyhead[L]{\small 基于 STM32F411 与 MAX30102 的脉搏测试仪}
\fancyhead[R]{\small 验收说明书}
\fancyfoot[C]{\thepage}
\renewcommand{\headrulewidth}{0.4pt}

\titleformat{\section}{\Large\bfseries\color{mainblue}}{\thesection}{0.8em}{}
\titleformat{\subsection}{\large\bfseries\color{darkgreen}}{\thesubsection}{0.8em}{}
\titleformat{\subsubsection}{\normalsize\bfseries}{\thesubsubsection}{0.7em}{}
\setlist[itemize]{leftmargin=1.8em,itemsep=2pt,topsep=3pt}
\setlist[enumerate]{leftmargin=1.8em,itemsep=2pt,topsep=3pt}
\renewcommand{\arraystretch}{1.25}
\setlength{\emergencystretch}{3em}
\sloppy
\newcolumntype{L}[1]{>{\raggedright\arraybackslash}p{#1}}

\lstdefinestyle{cstyle}{
  basicstyle=\ttfamily\scriptsize,
  backgroundcolor=\color{codebg},
  frame=single,
  rulecolor=\color{rulegray},
  breaklines=true,
  breakatwhitespace=false,
  columns=fullflexible,
  keepspaces=true,
  showstringspaces=false,
  xleftmargin=0.3em,
  xrightmargin=0.3em
}
\lstset{style=cstyle}

\newcommand{\keybox}[2]{
  \vspace{0.4em}
  \noindent
  \colorbox{#1}{%
    \parbox{\dimexpr\linewidth-2\fboxsep\relax}{#2}%
  }
  \vspace{0.4em}
}

\begin{document}

\begin{titlepage}
  \centering
  \vspace*{22mm}
  {\Huge\bfseries\color{mainblue} 基于 STM32F411 与 MAX30102 的脉搏测试仪\par}
  \vspace{8mm}
  {\LARGE 项目验收说明书（定稿版）\par}
  \vspace{12mm}
  {\large 心率、血氧、房颤风险提示与压力趋势指数的一体化实现\par}
  \vfill
  \begin{tabular}{rl}
    项目版本： & 2026-06-23 定稿版 \\
    核心源码： & \texttt{Core/Src/main.c} \\
    验收固件： & \texttt{build/Debug/STM32\_F411\_Test.elf} \\
    输出文档： & \texttt{output/pdf/验收说明\_脉搏测试仪\_2026-06-23.pdf} \\
  \end{tabular}
  \vfill
  {\large 项目组\par}
  {\large 2026-06-23\par}
\end{titlepage}

\tableofcontents
\newpage

\section{验收结论与功能边界}

本项目完成了基于 STM32F411 和 MAX30102 的多指标脉搏测试仪。设备通过 MAX30102 采集红光和红外 PPG 信号，在 STM32 端实时完成心率、血氧、房颤风险提示、压力趋势指数、OLED 显示、按键切换和 USB CDC 串口输出。

\keybox{softblue}{
\textbf{验收表述建议：}本项目实现的是课程设计级健康趋势监测原型。AF 页面显示的是基于 PPG/PPI 不规则性的房颤风险提示；Stress 页面显示的是基于 HRV 的压力趋势指数。二者均不是医学诊断结论。
}

\begin{longtable}{@{}L{3.2cm}L{10.8cm}@{}}
\toprule
项目 & 定稿内容 \\
\midrule
MCU & STM32F411 \\
传感器 & MAX30102，红光 + 红外双通道，100 Hz 采样 \\
显示 & 0.96 寸 OLED 显示心率、血氧、AF、Stress、AF test；0.91 寸 OLED 显示 PPG 波形与可靠峰标记 \\
心率 & IR 滤波窗口 + 自相关主周期估计 + 候选确认 + 跳变抑制 \\
血氧 & 红光/红外 AC/DC 比值法，公式为 \(SpO2=101-7R\) \\
AF 风险 & 11 维 PPI/RR 不规则性特征 + 离散朴素贝叶斯 \\
压力指数 & 13 维 HRV 特征 + 标准化逻辑回归 + 1--99 指数标定 \\
AF test & 内置公开 AFDB 06995 的 9030s--9060s AF rhythm RR/PPI 段，实时计算约 85\% 风险，不是硬编码概率 \\
\bottomrule
\end{longtable}

\section{项目文件夹结构}

下面只列验收需要重点说明的自主代码、数据、模型和报告。CubeMX/HAL 自动生成文件不展开。

\begin{lstlisting}
STM32_F411_Test/
|-- Core/Src/main.c
|   固件主程序：采样、滤波、心率、血氧、可靠峰、AF、Stress、OLED、按键、USB CDC
|-- USB_DEVICE/
|   STM32 USB CDC 虚拟串口
|-- build/Debug/STM32_F411_Test.elf
|   当前验收烧录固件
|-- scripts/read_max30102_raw.py
|   PC 端采集/回放脚本，能保存 raw.csv、processed.csv、summary.json
|-- scripts/train_af_naive_bayes.py
|   AF 公开 RR/PPI 数据处理、特征提取、朴素贝叶斯训练、JSON/C 导出
|-- scripts/validate_mcu_af_live.py
|   PC 端复刻 MCU AF live/test 策略，验证本地 PPG、公开数据和内置 test 段
|-- scripts/train_stress_hrv_model.py
|   WESAD BVP 处理、HRV 特征、逻辑回归训练、指数标定、JSON/C 导出
|-- scripts/validate_stress_hrv_live.py
|   本地 PPG 压力指数回放与质量门控验证
|-- training_dataset/
|   公开训练数据、派生窗口、模型参数、训练报告、全量验证报告
|-- testing dataset/
|   本地 MAX30102 采集 CSV、回放 CSV、故障复现和硬件验证数据
|-- AI助手接手？先读我.md
|   项目总入口和下一位接手者阅读顺序
|-- Reports/
|   进度归档、AF/压力情绪功能详细总结和补充汇报文档
|   |-- 总体完成进度（2026.6.23定稿）.md
|   |-- 房颤风险功能详细总结.md
|   |-- 压力情绪功能详细总结.md
|   |-- 脉搏测试仪项目补充汇报.docx
\end{lstlisting}

\section{整体数据与算法流水线}

\subsection{实时链路}

\begin{lstlisting}
MAX30102 红光/红外采样
  -> DC/AC 分离、RMS、手指质量判断
  -> 心率主支路：IR 滤波窗口 -> 自相关 -> 候选确认 -> HR 显示
  -> 血氧主支路：红光/红外 AC/DC 比值 -> SpO2 显示
  -> HRV 旁路：可靠峰确认 -> PPI 质量过滤
        -> 11 维 PPI 特征 -> AF 朴素贝叶斯 -> AF risk
        -> 13 维 HRV 特征 -> 逻辑回归 -> Stress index
  -> OLED/USB CDC 输出
\end{lstlisting}

设计要点是：心率/血氧主支路保持稳定显示，AF/压力作为旁路使用可靠峰 PPI。AF 和压力需要逐拍间期，因此不能直接用自相关后的平均心率代替。

\subsection{AF 与压力使用的生理评估指标}

\begin{longtable}{@{}L{2.2cm}L{3.1cm}L{3.5cm}L{5.7cm}@{}}
\toprule
功能 & 直接测量量 & 生理指标 & 本项目如何使用 \\
\midrule
AF 风险 & PPG 可靠峰得到的 PPI；公开训练中使用 ECG 标注得到的 RR 间期 & 心搏间期不规则性、脉搏节律离散程度、相邻间期突变比例 & 房颤常表现为逐拍间期明显不规则。模型提取平均 PPI、四分位距、变异系数、相邻 PPI 绝对差、pNN50/pNN80/pNN120 等 11 维特征，输出 AF-like 风险百分比。 \\
压力指数 & PPG/BVP 峰间期序列 & HRV，即心率变异性，反映自主神经调节状态 & 心理压力或高唤醒状态会影响交感/副交感平衡，使 HRV 指标变化。模型提取平均心率、SDNN、RMSSD、SDSD、pNN20/pNN50、CV、异常间期比例等 13 维特征，输出 1--99 压力趋势指数。 \\
\bottomrule
\end{longtable}

\section{心率与血氧算法}

\subsection{心率}

心率使用 IR 通道滤波后的窗口做自相关，寻找一段时间内的主周期。采样周期为 10 ms，如果自相关峰位于 \(k^\ast\)，则心率为：

\[
HR=\frac{60000}{10k^\ast}=\frac{6000}{k^\ast}
\]

为了避免瞬时噪声造成跳变，固件还加入候选确认、最近值中位数、跳变门限和超时恢复。主要实现位置：

\begin{itemize}
  \item 自相关和候选心率：\texttt{Core/Src/main.c:4510-4671}
  \item 心率确认和跳变保持：\texttt{Core/Src/main.c:4679-4912}
\end{itemize}

\subsection{血氧}

血氧使用红光与红外通道的 AC/DC 比值：

\[
R=\frac{AC_{red}/DC_{red}}{AC_{ir}/DC_{ir}},\qquad SpO2=101-7R
\]

固件只在手指质量、RMS、ratio 范围满足条件时更新显示。主要实现位置：\texttt{Core/Src/main.c:5028-5069}。

\section{AF 风险模型}

\subsection{数据来源}

AF 训练和验证使用 PhysioNet 公开节律数据，主要包括：

\begin{itemize}
  \item AFDB：Atrial Fibrillation Database；
  \item LTAFDB：Long-Term AF Database；
  \item NSR2DB：Normal Sinus Rhythm RR Interval Database。
\end{itemize}

本项目下载的是节律标注与 RR/PPI 间期相关文件，不依赖 ECG 大体积原始波形。这样做的原因是单片机实际只有 PPG，板端能稳定获得的是 PPI 间期，而不是 ECG 的 P 波、QRS 形态。

\subsection{11 维特征}

\begin{longtable}{@{}L{4.4cm}L{9.2cm}@{}}
\toprule
特征 & 含义 \\
\midrule
\path{valid_interval_count} & 窗口有效 PPI 数量 \\
\texttt{mean\_ppi\_ms} & 平均脉搏间期 \\
\texttt{iqr\_ppi\_ms} & PPI 四分位距 \(Q_{75}-Q_{25}\) \\
\path{trimmed_std_ppi_ms} & 去掉两端各 10\% 后的 PPI 标准差 \\
\texttt{trimmed\_cv\_ppi} & 截尾标准差 / 截尾均值 \\
\path{median_abs_delta_ms} & 相邻 PPI 绝对差中位数 \\
\path{p80_abs_delta_ms} / \path{p95_abs_delta_ms} & 相邻 PPI 绝对差的 80/95 分位数 \\
\path{pnn50_pct} / \path{pnn80_pct} / \path{pnn120_pct} & 相邻 PPI 差大于 50/80/120 ms 的比例 \\
\bottomrule
\end{longtable}

相邻差定义为：

\[
\Delta_i=PPI_i-PPI_{i-1},\qquad
pNN50=\frac{\#\{|\Delta_i|>50\text{ ms}\}}{N-1}\times100\%
\]

特征提取位置：\texttt{train\_af\_naive\_bayes.py:263-311}；MCU 等价实现位置：\texttt{Core/Src/main.c:1190-1304}。

\subsection{朴素贝叶斯实现原理}

训练时先把每个连续特征按固定边界离散成有限分箱，然后统计正常节律和 AF 节律在每个分箱内出现的概率。使用拉普拉斯平滑：

\[
P(x_i=b\mid c)=\frac{N_{c,i,b}+\alpha}{N_c+\alpha B_i},\qquad \alpha=1
\]

推理时使用对数概率累加，避免许多小概率相乘导致下溢：

\[
S_c=\log P(c)+\sum_{i=1}^{11}\log P(x_i\mid c)
\]

最后将 AF 类和 normal 类的分数差转换为风险百分比：

\[
P(AF\mid x)=\frac{1}{1+\exp[-(S_{AF}-S_{normal})]},\qquad
Risk_{AF}=100P(AF\mid x)
\]

\keybox{softgreen}{
\textbf{为什么可以用朴素贝叶斯：}AF 风险的主要线索是 PPI/RR 间期不规则性；这些特征可以离散为稳定分箱。模型训练后只需保存先验概率、分箱边界和条件概率表，MCU 端只做分箱、查表、累加和一次指数运算，计算量小、可解释、易验收。
}

代码片段（\texttt{Core/Src/main.c:1307-1355}）：

\begin{lstlisting}[language=C]
normal_score += af_nb_log_prob_normal[i][bin];
af_score     += af_nb_log_prob_af[i][bin];
diff = af_score - normal_score;
risk = 100.0f / (1.0f + expf(-diff));
\end{lstlisting}

\subsection{训练与验证结果}

\begin{longtable}{@{}L{4.7cm}L{3.6cm}L{5.3cm}@{}}
\toprule
数据范围 & 关键数量/指标 & 结果说明 \\
\midrule
公开数据总窗口 & 1,099,825 个窗口 & AF 窗口 320,784，normal 窗口 779,041 \\
公开数据总体 AUC & 0.969 & 表明 PPI/RR 不规则性特征对 AF 与 normal 有较强区分能力 \\
阈值 50\% 下召回率 & 0.9622 & AF 段大多数能被识别出来 \\
阈值 50\% 下特异度 & 0.9021 & normal 段大多数能保持低风险 \\
\bottomrule
\end{longtable}

报告文件：\texttt{training\_dataset/reports/房颤模型训练验证结果.json}。

\subsection{AF test 内置数据说明}

\keybox{softblue}{
\textbf{这一段是验收时最重要的解释：}AF test 显示约 85\% 不是写死的数值。单片机内部存的是一组公开 AFDB 记录的 RR/PPI 间期，进入 test 页面后按真实间期逐个播放，窗口满约 30 s 后调用同一套朴素贝叶斯模型实时算出约 85\%。
}

内置数据位于 \texttt{Core/Src/main.c:383-389}：

\begin{lstlisting}[language=C]
static const uint16_t af_test_ppi_ms[] = {
  /* AFDB 06995, 9030s..9060s, AF rhythm, 30s RR/PPI window, PC risk ~=85%. */
  540U, 800U, 784U, 556U, 800U, 760U, 780U, 796U, 792U, 752U, 756U, 776U,
  784U, 788U, 788U, 756U, 772U, 696U, 868U, 764U, 536U, 836U, 804U, 816U,
  804U, 784U, 772U, 800U, 808U, 404U, 1104U, 644U, 552U, 840U, 692U, 840U,
  764U, 760U, 776U, 800U
};
\end{lstlisting}

这组数据的来源和意义：

\begin{itemize}
  \item 来源是公开 PhysioNet AFDB 记录 \texttt{06995}；
  \item 时间段为 9030s--9060s；
  \item 标注节律为 AF rhythm；
  \item 数组中保存的是 40 个 RR/PPI 间期，单位 ms；
  \item 不是保存风险值，也不是保存显示文本；
  \item 固件播放逻辑位于 \texttt{Core/Src/main.c:848-920}；
  \item 每加入一个 PPI 后调用 \texttt{App\_UpdateAfTestRisk()}，窗口不足时显示等待，窗口满足后调用 \texttt{AF\_RiskPercentFromPpi()}；
  \item PC 等价验证结果为：\texttt{mcu\_test\_afdb\_06995}，首次有效时间 30.244 s，风险 85\%。
\end{itemize}

因此，AF test 的 85\% 是从公开 AF 节律间期、同一套 11 维特征和同一套朴素贝叶斯参数计算得到的结果，可复现、可追溯，不是伪造固定值。

\section{压力指数模型}

\subsection{数据来源与任务定义}

压力模型使用 WESAD（Wearable Stress and Affect Detection）公开数据集。训练信号选用腕部 BVP，采样率 64 Hz。BVP 与本项目 PPG 一样都是光电脉搏波，适合提取 PPI 和 HRV。

本项目不做开心、难过等具体情绪分类，只输出压力趋势指数：

\begin{longtable}{@{}L{3.0cm}L{4.0cm}L{6.5cm}@{}}
\toprule
指数范围 & OLED 状态 & 解释 \\
\midrule
1--29 & \texttt{relax} & 偏放松 \\
30--59 & \texttt{normal} & 日常正常区间 \\
60--79 & \texttt{medium} & 中等压力/唤醒水平 \\
80--99 & \texttt{high} & 偏高压力/唤醒水平 \\
\bottomrule
\end{longtable}

\subsection{13 维 HRV 特征}

\begin{longtable}{@{}L{4.1cm}L{9.5cm}@{}}
\toprule
特征 & 含义 \\
\midrule
\path{valid_interval_count} & 有效 PPI 数量 \\
\texttt{mean\_hr\_bpm} & 平均心率，\(60000/\overline{PPI}\) \\
\texttt{mean\_ppi\_ms} & 平均 PPI \\
\texttt{sdnn\_ms} & PPI 标准差，反映总体变异性 \\
\texttt{rmssd\_ms} & 相邻差平方均值平方根，常用于短时 HRV \\
\texttt{sdsd\_ms} & 相邻差标准差 \\
\path{pnn20_pct} / \path{pnn50_pct} & 相邻差大于 20/50 ms 的比例 \\
\texttt{cv\_ppi} & \(SDNN/\overline{PPI}\) \\
\path{median_abs_delta_ms} / \path{p80_abs_delta_ms} / \path{p95_abs_delta_ms} & 相邻差绝对值分布 \\
\texttt{outlier\_frac} & 相对中位数过短或过长的异常间期比例 \\
\bottomrule
\end{longtable}

特征提取位置：\texttt{train\_stress\_hrv\_model.py:188-223}；MCU 等价实现位置：\texttt{Core/Src/main.c:1597-1713}。

\subsection{逻辑回归实现原理}

每个 HRV 特征先按训练集均值和标准差标准化：

\[
\tilde{x}_i=\frac{x_i-\mu_i}{\sigma_i}
\]

然后计算线性分数和 stress 概率：

\[
z=b+\sum_{i=1}^{13}w_i\tilde{x}_i,\qquad
p=\frac{1}{1+e^{-z}}
\]

最后把概率按训练集 nonstress/stress 中位数标定到 1--99 指数。当前模型参数来自 \texttt{training\_dataset/models/压力HRV模型参数.json}：

\begin{itemize}
  \item 模型类型：\texttt{logistic\_regression}
  \item 类别：\texttt{nonstress} / \texttt{stress}
  \item 截距：\(-1.48795\)
  \item 概率标定：nonstress 中位数 0.06719 对应指数 45，stress 中位数 0.89503 对应指数 82
  \item 显示偏置：\(+5\)
\end{itemize}

\keybox{softgreen}{
\textbf{为什么可以用逻辑回归：}HRV 特征维度低、含义清楚，WESAD 样本规模有限。逻辑回归参数少、可解释，输出天然是 0--1 概率，便于映射到手表类似的 1--99 压力指数。移植到 MCU 时只需要保存均值、尺度、权重、截距和标定参数。
}

代码片段（\texttt{Core/Src/main.c:1716-1745}）：

\begin{lstlisting}[language=C]
scaled = (features[i] - stress_scaler_mean[i]) / stress_scaler_scale[i];
linear += stress_lr_coef[i] * scaled;
prob = 1.0f / (1.0f + expf(-linear));
index = Stress_IndexFromProbability(prob);
\end{lstlisting}

\subsection{训练与验证结果}

\begin{longtable}{@{}L{4.7cm}L{3.7cm}L{5.2cm}@{}}
\toprule
数据范围 & 指标 & 结果说明 \\
\midrule
WESAD 总窗口 & 3,066 个窗口 & nonstress 2,177，stress 889 \\
WESAD 总体 AUC & 0.9412 & HRV 特征对压力/非压力有明显区分能力 \\
阈值 60 下召回率 & 0.9539 & stress 窗口大多能被识别为中高压力 \\
baseline 中位指数 & 50 & 与用户手表日常约 50 左右的读数接近 \\
stress 中位指数 & 87 & 公开压力任务下明显升高 \\
\bottomrule
\end{longtable}

报告文件：\texttt{training\_dataset/reports/压力HRV模型训练验证结果.json}。

高心率场景下，运动造成的心率升高和 HRV 改变也可能推高压力指数。最终固件在 \(HR>120\) bpm 时仍后台计算，但 OLED 显示 \texttt{-- high HR}，避免把运动反应误解释为心理压力。

\section{MCU 端移植与显示策略}

\subsection{关键参数}

\begin{longtable}{@{}L{4.8cm}L{2.8cm}L{6.4cm}@{}}
\toprule
参数 & 当前值 & 位置/含义 \\
\midrule
AF 特征数 & 11 & \texttt{Core/Src/main.c:264} \\
AF 最少 PPI & 20 & \texttt{Core/Src/main.c:269} \\
AF 窗口 & 30 s & \texttt{Core/Src/main.c:270} \\
AF 快/慢更新 & 10 s / 30 s & 高风险或未出值阶段快更新，低风险稳定后慢更新 \\
AF 最大向上跳变 & 20 个百分点 & 避免正常测量瞬间跳到 90--100\% \\
Stress 特征数 & 13 & \texttt{Core/Src/main.c:279} \\
Stress 最少 PPI & 28 & \texttt{Core/Src/main.c:280} \\
Stress 窗口 & 40 s & 实机缩短首次显示等待 \\
Stress 首次/刷新 & 10 s / 30 s & 首次或 high HR 状态更快尝试，稳定后 30 s 刷新 \\
Stress 高心率隐藏 & HR \(>120\) bpm & OLED 显示 \texttt{-- high HR} \\
\bottomrule
\end{longtable}

\subsection{按键页面}

当前按键页面顺序：

\begin{enumerate}
  \item 心率/血氧主页面；
  \item AF Live：实时 AF risk；
  \item Stress：第一行显示数值，第二行显示 \texttt{relax/normal/medium/high} 或 \texttt{-- high HR}；
  \item AF Test：播放内置 AFDB 06995 PPI 段，展示完整 AF 推理链路。
\end{enumerate}

\subsection{质量门控}

AF 和 Stress 共享可靠峰 PPI 质量门控：

\begin{itemize}
  \item PPI 必须处于合理生理范围；
  \item PPI 要与当前自相关心率大体一致；
  \item 当前窗口拒绝比例过高时不更新，保留旧值；
  \item AF 使用上跳门限，避免单个异常窗口造成突然高风险；
  \item Stress 在高心率时隐藏前台数值，但后台仍计算。
\end{itemize}

PPI 过滤和窗口质量实现位置：\texttt{Core/Src/main.c:4339-4425}。

\section{PC 脚本与模型文件说明}

\begin{itemize}
  \item \path{scripts/read_max30102_raw.py}：串口采集与离线回放 MAX30102 原始数据，输出 raw/processed/summary，辅助定位硬件、峰检测和显示问题。
  \item \path{scripts/train_af_naive_bayes.py}：从 PhysioNet RR/PPI 标注构造窗口，提取 11 维 AF 特征，训练离散朴素贝叶斯，导出 JSON 和 C 参数。
  \item \path{scripts/validate_mcu_af_live.py}：复刻 MCU AF 策略，验证 live AF、内置 AFDB test 段、公开数据、本地 PPG 和跳变保持逻辑。
  \item \path{scripts/train_stress_hrv_model.py}：从 WESAD BVP 提取 PPI/HRV，训练标准化逻辑回归，完成 1--99 指数标定。
  \item \path{scripts/validate_stress_hrv_live.py}：用本地 PPG 回放验证 Stress 指数、质量门控和 high HR 显示策略。
  \item \path{training_dataset/models/}：保存 AF 朴素贝叶斯参数、Stress 逻辑回归参数，以及导出到 MCU 的 C 头文件。
  \item \path{training_dataset/reports/}：保存公开数据训练指标、本地回放、全量验证和当前定稿报告。
  \item \path{testing dataset/}：保存本地采集 CSV，包括 raw.csv、processed.csv、events.csv、summary.json 和 replay 文件。
\end{itemize}

\section{烧录与验收演示}

\subsection{DAPLink 烧录}

硬件连接：

\begin{itemize}
  \item DAPLink GND 接开发板 GND；
  \item DAPLink CK/TCK 接 SWCLK；
  \item DAPLink IO/TMS 接 SWDIO；
  \item 若开发板由 Type-C 供电，DAPLink 3V3 不接。
\end{itemize}

编译：

\begin{lstlisting}
& '%LOCALAPPDATA%\stm32cube\bundles\cmake\4.3.1+st.1\bin\cmake.exe' --build build\Debug
\end{lstlisting}

烧录：

\begin{lstlisting}
& 'tools\xpack-openocd-0.12.0-7\bin\openocd.exe' `
  -s 'tools\xpack-openocd-0.12.0-7\openocd\scripts' `
  -f interface\cmsis-dap.cfg -f target\stm32f4x.cfg `
  -c 'adapter speed 1000; program build/Debug/STM32_F411_Test.elf verify reset exit'
\end{lstlisting}

成功标志：\texttt{Programming Finished}、\texttt{Verified OK}、\texttt{Resetting Target}。

\subsection{推荐演示顺序}

\begin{enumerate}
  \item 进入心率/血氧页面，手指稳定后说明 HR 与 SpO2 来自主支路。
  \item 切到 AF Live，说明需要约 30 s PPI 窗口，正常人通常低风险。
  \item 切到 Stress，说明需要约 40 s PPI 窗口，稳定后 30 s 刷新。
  \item 切到 AF Test，说明内置的是公开 AFDB 06995 的 AF rhythm 间期，不是固定 85\%，等待约 30 s 后展示计算结果。
  \item 如老师问数据真实性，直接展示 \texttt{Core/Src/main.c:383-389} 的数组注释和 \texttt{validate\_mcu\_af\_live.py} 的等价验证。
\end{enumerate}

\section{答辩常见问题}

\subsection{AF 是人工智能吗？}

是机器学习分类模型，不是神经网络。AF 使用公开 RR/PPI 数据训练离散朴素贝叶斯，板端执行概率推理。

\subsection{为什么不用 ECG？}

本设备硬件只有 MAX30102，能测 PPG，不具备 ECG 形态信息。因此 AF 只能做基于 PPI 不规则性的风险提示，不能做临床房颤诊断。

\subsection{为什么 AF test 是 85\%，是不是写死的？}

不是。固件内置的是公开 AFDB 06995 的一段 AF rhythm RR/PPI 间期。test 页面逐个播放这些间期，窗口满足后调用同一套 AF 特征和朴素贝叶斯推理，PC 等价验证得到约 85\%。

\subsection{压力为什么不是开心/难过分类？}

设备只有 PPG，没有 EDA、呼吸、皮温、语音或上下文。仅靠 HRV 更适合做压力/唤醒程度趋势，不适合可靠区分具体情绪。

\subsection{这个项目最终完成了什么？}

最终完成了一个可烧录、可演示、可复现的数据闭环：传感器采集、心率血氧计算、可靠峰和 PPI 质量控制、AF 机器学习风险提示、压力 HRV 机器学习指数、OLED 显示、串口采集、PC 训练验证和验收文档。

\end{document}
